\documentclass{SWUSTart}
\usepackage[utf8]{inputenc}
\usepackage[T1]{fontenc}
\usepackage{geometry}
\usepackage{graphicx}
\usepackage{float}
\usepackage{hyperref}
\usepackage[english,indonesian]{babel}

\title{Optimasi Pemrosesan Template LaTeX Menggunakan Pendekatan Berbasis AI}
\author{Academic Support Agent}
\date{\today}

\begin{document}
\maketitle

\begin{abstract}

\end{abstract}

\chapter{PENDAHULUAN}


\chapter{TINJAUAN PUSTAKA}


% Original content (commented out):
% %!TEX encoding = UTF-8 Unicode
% !TeX program = xelatex
\documentclass[a4paper,bachelor,12pt]{SWUSTart}
%类别选项：bachelor, CDIOC,CDIOD,CDIOI,CDIOO, design,courseproject 
\renewcommand{\baselinestretch}{1.4}%重新定义行距
\begin{document}
	\definecolor{backgroundcolor}{RGB}{202,234,206}
	\pagecolor{backgroundcolor}%背景护眼色	
	\maketitle
%%%%%%%%%%%%%%	
\pagestyle{fancy}
%\linenumbers %行号
\tableofcontents  %目录	
% 表格索引
\listoftables
% 插图索引
\listoffigures	
\mainmatter%正文用页码阿拉伯数字
\setcounter{page}{1}%页码重新编号
\include{chapter/guider}
\include{chapter/chap2}
\include{chapter/chap1.tex}	
\include{chapter/chap3}
\backmatter %正文后页码用大写罗马数字
\include{chapter/conclusion}
\include{chapter/acknol}%致谢
\bibliography{chapter/standard} % 
\addcontentsline{toc}{chapter}{参考文献}
\appendix %开始附录
\include{chapter/researchresults}% 研究成果
\end{document}


\end{document}